% Options for packages loaded elsewhere
\PassOptionsToPackage{unicode}{hyperref}
\PassOptionsToPackage{hyphens}{url}
\PassOptionsToPackage{space}{xeCJK}
%
\documentclass[
]{article}
\usepackage{amsmath,amssymb}
% \usepackage{lmodern}
\usepackage{iftex}
\ifPDFTeX
  \usepackage[T1]{fontenc}
  \usepackage[utf8]{inputenc}
  \usepackage{textcomp} % provide euro and other symbols
\else % if luatex or xetex
  \usepackage{unicode-math}
  \defaultfontfeatures{Scale=MatchLowercase}
  \defaultfontfeatures[\rmfamily]{Ligatures=TeX,Scale=1}
  \setmainfont[]{DejaVu Serif}
  \setmathfont[]{Asana Math}
  \ifXeTeX
    \usepackage{xeCJK}
    \setCJKmainfont[]{Noto Serif CJK TC}
  \fi
  \ifLuaTeX
    \usepackage[]{luatexja-fontspec}
    \setmainjfont[]{Noto Serif CJK TC}
  \fi
\fi
% Use upquote if available, for straight quotes in verbatim environments
\IfFileExists{upquote.sty}{\usepackage{upquote}}{}
\IfFileExists{microtype.sty}{% use microtype if available
  \usepackage[]{microtype}
  \UseMicrotypeSet[protrusion]{basicmath} % disable protrusion for tt fonts
}{}
\makeatletter
\@ifundefined{KOMAClassName}{% if non-KOMA class
  \IfFileExists{parskip.sty}{%
    \usepackage{parskip}
  }{% else
    \setlength{\parindent}{0pt}
    \setlength{\parskip}{6pt plus 2pt minus 1pt}}
}{% if KOMA class
  \KOMAoptions{parskip=half}}
\makeatother
\usepackage{xcolor}
\IfFileExists{xurl.sty}{\usepackage{xurl}}{} % add URL line breaks if available
\IfFileExists{bookmark.sty}{\usepackage{bookmark}}{\usepackage{hyperref}}
\hypersetup{
  hidelinks,
  pdfcreator={LaTeX via pandoc}}
\urlstyle{same} % disable monospaced font for URLs
\setlength{\emergencystretch}{3em} % prevent overfull lines
\providecommand{\tightlist}{%
  \setlength{\itemsep}{0pt}\setlength{\parskip}{0pt}}
\setcounter{secnumdepth}{-\maxdimen} % remove section numbering
\usepackage[margin=1in]{geometry}
\usepackage{microtype}
\usepackage{hyperref}
\hypersetup{colorlinks=true, linkcolor=blue, urlcolor=blue}

\usepackage{fontspec}
\usepackage{xeCJK}

\setmainfont{DejaVu Serif}
\setCJKmainfont{Noto Serif CJK TC}[BoldFont = Noto Serif CJK TC]

\usepackage{amsmath, amssymb, amsthm}
\ifLuaTeX
  \usepackage{selnolig}  % disable illegal ligatures
\fi

\author{}
\date{}

\begin{document}

昨天我們討論了 Hamming Code
，是一個可以非常快速糾正錯誤的編碼系統。今天我們來看一個特別的「多項式編碼」：Reed-Solomon
Code

RS
編碼的核心思想是將訊息視為一個多項式，然後在有限域（你可先想成我們之前做的整數商環
Z\_q）上對其進行運算。這種方法使得 RS
編碼能夠高效地處理錯誤，並具有靈活的錯誤更正能力。今天我們先討論編碼的部分，明天再來看糾錯機制

\hypertarget{rs-ux7de8ux78bcux7cfbux7d71}{%
\section{RS 編碼系統}\label{rs-ux7de8ux78bcux7cfbux7d71}}

首先要選擇參數 k : 訊息的位元數 n : 編碼的位元數 q : 整數商環 Z\_q
的模數 其中需滿足 \[
k\leq n\leq q
\] 除此之外還需在 Z\_q 中選擇 n 個相異的點 \[
(\alpha_{1},\dots,\alpha_{n})
\] 你可以把這些點叫做「插值點」（或者「評估點」）

選擇要傳送的訊息 \[
m = (m_{0},m_{1},\dots,m_{k-1}).
\] 其中 m\_i 都是 Z\_q 中的點（可以相同）。

編碼方式如下： 生成多項式 \[
f_{m}(x) = m_{0}+m_{1}x+m_{2}x^{2}+\dots+m_{k-1}x^{k-1}
\] 你可以把它叫做「訊息多項式」

然後計算所有 alpha\_i 代進去的值 \[
(f_{m}(\alpha_{1}),f_{m}(\alpha_{2}),\dots,f_{m}(\alpha_{n}))
\] 這就完成了 RS 編碼！

\hypertarget{ux751fux6210ux77e9ux9663}{%
\subsection{生成矩陣}\label{ux751fux6210ux77e9ux9663}}

好的，我們進行以下計算 \[
(f_{m}(\alpha_{1}),\dots,f_{m}(\alpha_{n}))
= 
\left( \sum_{i=0}^{k-1} m_{i}\alpha_{1}^{i},
\sum_{i=0}^{k-1} m_{i}\alpha_{2}^{i},
\dots,
\sum_{i=0}^{k-1} m_{i}\alpha_{n}^{i} \right)
\] 於是可以把碼字寫成 mG ，其中 \[
G = 
\begin{bmatrix}
1  & 1 & \cdots & 1 \\
\alpha_{1} & \alpha_{2} & \dots & \alpha_{n} \\
\alpha_{1}^{2} & \alpha_{2}^{2} & \dots & \alpha_{n}^{2} \\
 \vdots & \vdots & \ddots & \vdots \\
\alpha_{1}^{k-1} & \alpha_{2}^{k-1} & \dots & \alpha_{n}^{k-1} \\
\end{bmatrix}
\]

\hypertarget{ux95dcux65bcux89e3ux78bc}{%
\subsection{關於解碼}\label{ux95dcux65bcux89e3ux78bc}}

假設傳遞過程沒有錯誤，從碼字 \[
(f_{m}(\alpha_{1}),f_{m}(\alpha_{2}),\dots,f_{m}(\alpha_{n}))
\] 解回 m
就是標準的多項式插值法，最常見的方法莫屬拉格朗日插值法，若你不怕計算麻煩，使用牛頓插值法或任何其他插值法也都可以。

但是當傳遞過程中出現錯誤時，假設最多有 t 個錯誤且收到的字為 \[
c = (c_{1},\dots,c_{n})
\] 那我們就面臨「在 c 中有最多其中 t
個是錯誤資料其他都是正常的函數值，要找到本來的 k-1 次多項式 f\_m
的問題」。乍看之下非常複雜，實際處理起來也是有點麻煩，我們明天再來研究這個糾錯機制。

\hypertarget{sagemath-ux5be6ux4f5c}{%
\section{SageMath 實作}\label{sagemath-ux5be6ux4f5c}}

\begin{verbatim}
# Parameters:

q = 17
n = 7
k = 3

R = quotient(ZZ, q*ZZ)

alphas = []
while len(alphas) < n:
    elem = R.random_element()
    if elem not in alphas:
        alphas.append(elem)

print(f"插值點: {alphas}")

# Outputs:
# 插值點: [13, 16, 7, 14, 2, 9, 1]
\end{verbatim}

\begin{verbatim}
# Generator matrix

G = Matrix(
    R,
    [
        [alphas[i]^j for i in range(n)]
        for j in range(k)
    ]
)
print("生成矩陣 G:")
print(G)

# Ouptuts:
# 生成矩陣 G:
# [ 1  1  1  1  1  1  1]
# [13 16  7 14  2  9  1]
# [16  1 15  9  4 13  1]
\end{verbatim}

\begin{verbatim}
message = [R.random_element() for _ in range(k)]

print(f"訊息: {message}")

# Outputs:
# 訊息: [6, 3, 1]
\end{verbatim}

\begin{verbatim}
def encode(message):
    return (Matrix(message) * G).list()

encoded = encode(message)

print(f"編碼: {encoded}")

# Outputs:
# 編碼: [10, 4, 8, 6, 16, 12, 10]
\end{verbatim}

\hypertarget{takeaway}{%
\section{Takeaway}\label{takeaway}}

\begin{itemize}
\tightlist
\item
  多項式編碼很籠統地來說就是使用多項式這個數學物件所做的編碼，一個重要的例子就是
  Reed-Solomon Code
\item
  Reed-Solomon
  的編碼方式是將訊息作為係數來產生一個多項式，然後把插值點帶進去。
\end{itemize}

ref: GURUSWAMI, Venkatesan; RUDRA, Atri; SUDAN, Madhu. Essential coding
theory.

\end{document}
